% 1. Preamble
\documentclass[12pt, a4paper]{article}

% Import Packages
\usepackage[utf8]{inputenc}        % Encoding
\usepackage{amsmath, amssymb}     % Advanced math formulas
\usepackage{mathtools}           % Math enhancements
\usepackage{graphicx}            % Image insertion
\usepackage{hyperref}            % Interactive hyperlinks & PDF bookmarks
\usepackage{color}               % Color text and backgrounds
\usepackage{geometry}            % Page geometry
\usepackage{comment}             % Multi-line comments
\usepackage{titling}             % Custom title formatting

\geometry{a4paper, margin=1in}

% Custom Command Definitions
\newcommand{\Sball}{\mathcal{S}}
\providecommand{\norm}[1]{\left\lVert#1\right\rVert}

% Hyperref Setup for Clean PDF Links
\hypersetup{
    colorlinks=true,
    linkcolor=blue,
    urlcolor=blue,
    pdftitle={Pattern Arithmetic: Verdicts from Bowls and Paths},
    pdfauthor={Sandeep Lakshminarayan Chiluveru}
}

% Metadata & Custom Title Banner
\pretitle{\begin{center}\Huge\fontfamily{tt}\selectfont ================================================================================\\[0.5em]}
\posttitle{\\[0.5em] ================================================================================\end{center}}

\title{Pattern Arithmetic: Verdicts from Bowls and Paths}
\author{Sandeep Lakshminarayan Chiluveru}
\date{\today}

\begin{document}

\maketitle

% ABSTRACT
\begin{abstract}
A companion paper, \emph{Pattern Arithmetic: Typed Algebra of Bounded
Units on a Phase-Colored Sphere}, defines a unit of truth
$x=(r,\hat n,\chi,\tau)$ on a sphere $\mathcal{S}$ and an algebra over
such units --- gather ($\oplus$), remove ($\ominus$), mix ($\boxplus$),
walk ($T$), and offset --- together with an admissibility test that
decides which sorts $\tau$ may use which operators truthfully. That
paper builds the algebra and stops at the reading: a bowl of units may
be mixed or walked, and the result is a state. It does not say how a
state, or a bowl, becomes a \emph{decision}.

This paper adds the missing step. A shaman, a tribal meteorologist, or
a market analyst holds a handful of signs and must arrive at one of a
small number of conclusions --- water or no water, storm or calm, rise
or fall. We show that the algebra already contains two distinct
mechanisms for this, corresponding to the two readings the first paper
leaves undemonstrated: an order-free reading, in which evidence
accumulates as phasors and a balanced case returns no answer rather
than a forced one; and an order-dependent reading, in which evidence
is walked as a sequence of actors to a landing compared against
seated names. We give a rule --- the \emph{swap test} --- for which
reading a given body of evidence calls for, and identify four
components a decision needs beyond what the algebra supplies: a
comparison rule, a selection rule for which signs enter and in what
order, a principled abstention, and a place for a lexicon of
conclusions to be learned rather than declared. This paper is a
framework: it specifies these four components precisely enough to be
filled in, and works one example --- the diviner's shells --- by hand.
It does not report an evaluation against data, and says exactly where
that evaluation would have to be made.
\end{abstract}

\section{Introduction}
\label{sec:intro}

A shaman throws ten shells and reads a future from them. He may throw
twenty and read it better, or forty better still --- but unless he
says so, no one else knows what he saw. The shells are evidence; the
saying is the verdict. Pattern Arithmetic gives a precise account of
the shells. It has none, so far, of the saying.

This is deliberate on the first paper's part, not an oversight. That
paper's discipline was to make the algebra trustworthy before asking
it to decide anything: an admissibility test that refuses sorts with
no honest map onto the operators, a type table that says which
combinations are legal, and a running insistence that a walk's landing
is compared to a name, never confused with one. Asking that paper to
also commit to a decision rule would have asked it to defend a claim
it was not ready to make.

It is ready now. This paper takes the atom, the bowl, and the path as
given --- cited, not re-derived --- and asks the next question: given
a bowl of evidence, what is the answer?

\subsection{A human can hold a handful; a machine need not stop
there}
A diviner's working memory bounds how many shells, how many birds, how
many signs can be weighed against one question at once. Nothing in
the algebra of the first paper imposes that bound. The claim of this
paper is narrower than it might sound: not that more evidence is
always better, but that \emph{the mechanism for combining evidence
does not change shape as the amount of evidence grows} --- the same
gather, the same two readings, apply to ten shells or ten thousand
sensor readings. Where that mechanism strains under scale (\S 4.3) is
stated plainly rather than smoothed over.

\section{What Is Inherited}
\label{sec:inherited}

We use, without re-deriving, the following from the companion paper.
Equation and axiom numbers refer to that paper.

\begin{itemize}
\item The atom $x=(r,\hat n,\chi,\tau)\in\mathcal{S}$, with $r\le 1$ a
      normalized strength, $\hat n\in S^2$ a direction, $\chi\in[0,2\pi)$
      a phase, and $\tau$ a sort tag that gates which operators may be
      truthfully applied (Axiom 1; Axiom 0).
\item Gather, $\oplus$: union on finite subsets of $\mathcal{S}$,
      sort-blind, idempotent, lossless, undone by $\ominus$ (Axiom 4).
\item Mix, $\boxplus$: a phasor reading of a bowl whose members share a
      direction and a sort,
      $\boxplus(B)=\bigl(\min(1,|z_B|),\hat n,\arg z_B,\tau\bigr)$
      with $z_B=\sum_{x\in B} r_x e^{i\chi_x}$; commutative,
      non-associative, lossy, uninvertible (Axiom 5).
\item The walk, $T$: an ordered composition
      $T_{w_n}\circ\cdots\circ T_{w_1}(x)$ of actor-parameterized
      rotations and phase shifts; non-commutative in general, and
      many-to-one, so a landing alone does not determine the walk that
      reached it (Axiom 6; \S3).
\item \emph{Levels are closed}: a unit carries no representation of
      what it constitutes, and a bowl carries no representation of the
      reading that will be applied to it. Whatever a lower unit would
      need to know about a higher one --- which units were taken as
      co-present, over what interval, under what question --- is a
      constituting fact of the higher level, not a coordinate of the
      atom (companion paper, \S2.2).
\end{itemize}

Two consequences of level closure matter enough here to restate. First,
a bowl that saturates under Mix is often a level that has been
flattened --- repeated samples of one unit, mistaken for many units at
one level (companion paper, \S2.7). Second, the atom stays a
four-tuple; grouping criteria such as a time window or a spatial
radius live outside it, deciding bowl membership before the algebra
runs rather than riding inside any single unit.

\subsection{A substitution: bounding aggregated strength}
\label{subsec:bijection_substitution}

The companion paper's Axiom 1 names its clip, $r'=\min(1,|z|)$, as a
choice, and names an order-preserving bijection as the alternative a
setting that aggregates at scale should weigh --- naming this paper's
own use as exactly such a setting, rather than leaving that
identification for the reader to make. This paper takes that
invitation: wherever a strength is reported from a phasor sum for
reading purposes (Reading A's own verdict, and a reference atom's own
reported strength), it substitutes
\begin{equation}
    r' \;=\; \frac{|z|}{K+|z|},
    \label{eq:bijection}
\end{equation}
the $K$-calibrated case of the general Hill form $x^n/(K^n+x^n)$ at
$n=1$, for the companion paper's clip. $K$ is chosen per domain, not
fixed here: it sets where the half-confidence point sits on the scale
of that domain's own evidence. Held at $n=1$, this does not trade away
the reason a bijection was wanted in the first place --- the tail
behaviour $1-r'\approx K/|z|$ for large $|z|$ still lets arbitrarily
large aggregates be distinguished, only at a correspondingly large
scale set by $K$. Raising $n$ above $1$ is not free in the same way:
it steepens the transition near $K$ at the cost of exactly that tail,
reintroducing practical indistinguishability at a smaller scale rather
than removing it, and is not adopted here.

This substitution leaves $\bar R_c$ (Equation~\eqref{eq:coherence})
untouched: that quantity already normalizes by the raw, uncapped
$|z|$, and was built for the same reason this substitution is --- to
stop a bound from erasing a real difference in agreement.

\section{Two Readings of One Bowl}
\label{sec:two_readings}

Given a bowl $B=\{x_1,\ldots,x_n\}$ of evidence bearing on one
question, the companion paper currently defines two readings that turn
$B$ into a single state: phasor Mix and ordered transformation. This is
not a claim that no other reading could be defined; it is the pair
that paper supplies today, and the pair this paper uses. We name them
here because the first paper defines both and demonstrates only one.

\subsection{Reading B: the walk}
\label{subsec:reading_b}
Seat the question as a state $q$. Apply the members of $B$ as actors,
in some order $\pi$:
\[
    \mathrm{land}_\pi(q) \;=\; T_{x_{\pi(n)}}\circ\cdots\circ
    T_{x_{\pi(1)}}(q).
\]
The verdict is the landing, compared against seated names (the storm's
$\mathrm{calm}$ and $\mathrm{storm}$ pins are an instance of this).
Order matters by construction: $T_a T_b\neq T_b T_a$ in general, so
$\mathrm{land}_\pi \neq \mathrm{land}_{\pi'}$ for $\pi\neq\pi'$ is the
expected case, not a failure. This reading is fully worked in the
companion paper's storm example, which we do not repeat.

\subsection{Reading A: the phasor sum}
\label{subsec:reading_a}
Rotate every member of $B$ onto the question's axis $\hat n_q$ (a
rotation, which preserves $r$ and $\chi$), then read the whole bowl at
once:
\[
    \boxplus(B) \;=\; \Bigl(\min(1,|z_B|),\ \hat n_q,\ \arg z_B,\ \tau\Bigr),
    \qquad
    z_B=\sum_{x\in B} r_x e^{i\chi_x}.
\]
Order plays no role: $z_B$ is a sum, and sums do not care about the
order of their terms. The verdict's \emph{confidence} is $|z_B|$'s
image under the cap, and its \emph{direction} is $\arg z_B$, read
against a lexicon of conclusions seated at particular phases (e.g.\
$\chi=0$ for ``water,'' $\chi=\pi$ for ``no water''). Agreeing evidence
lengthens the sum; opposing evidence shortens it; evenly opposed
evidence returns $z_B=0$ --- black, and a principled abstention rather
than a forced choice (companion paper, Axiom 2).

\subsection{The swap test}
\label{subsec:swap_test}
Which reading applies to a given $B$ is decided by one question: does
exchanging the order in which two members of $B$ arrived change the
answer?
\begin{itemize}
\item If \textbf{no} --- the shells thrown together, ten readings
      taken at one instant --- $B$ has no order to respect, and
      Reading A is the correct mechanism. Applying Reading B here
      manufactures an order-dependence that is not in the evidence.
\item If \textbf{yes} --- birds fleeing, then a wind shift, then a
      temperature drop, where the same three signs in a different
      sequence mean something else --- $B$'s order is itself
      information, and Reading A would discard it. Reading B is
      required.
\end{itemize}
The two readings are not competing hypotheses about one mechanism;
they are the correct tool for two different shapes of evidence, and
the swap test is how a system --- human or otherwise --- tells which
shape it is holding.

\subsection{The readings compose}
\label{subsec:readings_compose}
Nothing prevents using both on one episode. A morning's observations
may arrive in a meaningful sequence (Reading B over the morning) and
also include several simultaneous signs at one moment within it
(Reading A over that moment's bowl, itself one actor in the larger
walk). This is level closure in practice (\S\ref{sec:inherited}): a
Reading A verdict is a state, and a state may act as $w$ in a $T$-path,
exactly as any other seated unit may.

\section{Four Components the Algebra Does Not Supply}
\label{sec:four_gaps}

Both readings above take a bowl and return a state. Neither, on its
own, returns a decision a person or a system can act on. Four
components stand between a state and a verdict, and none of them is
supplied by $\oplus$, $\boxplus$, $T$, or offset.

\subsection{A comparison rule}
\label{subsec:comparison_rule}
The companion paper's storm section says a ball falls \emph{near} the
$\mathrm{calm}$ pin or the $\mathrm{storm}$ pin, and does not define
near. A working comparison rule must state, at minimum: which
coordinates of the landing are read, what metric is used on those
coordinates, and whether $\tau$ must match exactly or may differ.

Two states are compared only if $\tau_A=\tau_B$; $\tau$ gates
comparison the same way it gates every other operator. Given that, a
distance is taken over all three geometric coordinates, each measured
the way the companion paper already measures it elsewhere rather than
by naive coordinate subtraction:
\begin{equation}
    d(A,B)^2 \;=\;
    w_r\,(r_A-r_B)^2
    \;+\;
    w_{\hat n}\,\arccos(\hat n_A\cdot\hat n_B)^2
    \;+\;
    w_\chi\,\bigl(\min(|\chi_A-\chi_B|,\,2\pi-|\chi_A-\chi_B|)\bigr)^2 ,
    \label{eq:comparison_metric}
\end{equation}
using the geodesic distance the offset operator already relies on for
$\hat n$, and the circular distance a phase requires for $\chi$; $r$
alone is a bounded line with no wraparound and needs neither
correction. Raw coordinate subtraction on $\hat n$ or $\chi$ is
\emph{not} a safe default: two directions or phases near $0$ and near
$2\pi$ are adjacent on the sphere or the circle but appear maximally
distant under naive subtraction, and reference atoms placed near a
natural zero point --- which is common, not rare --- would be
misjudged systematically rather than occasionally. \emph{Open in this
paper:} the weights $w_r,w_{\hat n},w_\chi$, meaning how much a
mismatch in strength, direction, or phase should count against an
otherwise close match. \S\ref{sec:worked_example} does not need this
choice made in general: every shell and every reference point there
already shares $\hat n_q$ by construction, so $w_{\hat n}$ never
enters, and the verdict is read from $\chi$ (which reference point is
nearest) with $r$ supplying the null rather than a weighted distance.
A question with reference atoms at different directions would need
$w_{\hat n}$ and $w_\chi$ actually chosen against each other, which
remains open.

\paragraph{A similarity metric.} $d(A,B)$ answers how far apart two
atoms are; the same question read the other way round --- how alike
--- has never been named, though the raw ingredient has been sitting
inside Equation~\eqref{eq:comparison_metric} since it was written.
Direction similarity is exactly the dot product the geodesic term
already computes, before $\arccos$ turns it into an angle:
\begin{equation}
    \mathrm{sim}(A,B) \;=\; \hat n_A\cdot\hat n_B.
    \label{eq:direction_similarity}
\end{equation}
Bounded $[-1,1]$ by construction --- $+1$ identical direction, $0$
orthogonal, $-1$ antipodal --- and exactly invertible against the
geodesic term already in use: $\arccos(\mathrm{sim}(A,B))$ recovers
that term precisely, confirmed to four decimal places across a range
of directions. This is deliberately the narrow, direction-only
reading, matching how cosine similarity is ordinarily used elsewhere
(comparing embedding directions, not a magnitude-weighted blend); it
needs none of $w_r,w_{\hat n},w_\chi$, so it is available today
regardless of how that open choice is eventually settled. A fuller
similarity built from the whole of $d(A,B)$ --- such as $e^{-d(A,B)}$
--- would answer a related but different question, weighting
strength and phase alongside direction, and would inherit the same
open weights $d(A,B)$ itself depends on; left for whichever setting
first needs that fuller reading rather than adopted here.

\subsection{A selection rule}
\label{subsec:selection_rule}
For Reading B, the order $\pi$ that the evidence is applied in is
handed to the walk; it is not the walk's to choose. A real system must
decide which signs enter $B$ at all (\S\ref{subsec:comparison_rule}'s
window and radius criteria, companion paper \S2.2) and, when order
matters, in what sequence. This selection is itself where most of an
inference system's actual difficulty lives, and it is entirely outside
the algebra. \emph{Open in this paper.}

\paragraph{A sequence is a degenerate graph.} A walked path
$T_{w_n}\circ\cdots\circ T_{w_1}(x)$ is a simple directed path, one
actor after another with no branching. When $\pi$ is not given in
advance, the natural structure is not a single path but a tree of
candidate paths: several possible next actors branch from a shared
prefix (twinkle, twinkle, then \emph{either} little or something
else), and choosing a walk means searching that tree rather than
replaying one. Every edge of the tree is still one application of
$T$; the tree organizes a search over sequences, it does not add an
operator. How candidates branch, how a branch is scored before it is
fully walked, and how the tree is pruned are the actual content of
this open selection rule, not resolved by naming the structure.

\subsection{A principled null}
\label{subsec:the_null}
Reading A supplies one for free: $z_B=0$. Reading B does not. A walk
always lands somewhere, and a comparison rule (\S\ref{subsec:comparison_rule})
always returns a nearest pin, however far away it is. A Reading-B
system therefore needs an explicit distance threshold beyond which the
verdict is \emph{abstain} rather than \emph{nearest pin regardless}.
This is a second place, distinct from \S\ref{subsec:comparison_rule},
where a number must be chosen rather than derived.

\subsection{Learned pins and learned signs}
\label{subsec:learned_pins}
The companion paper notes that a lookup $E_L$ may be learned
(\S5.6) without the units it points into being anything other than
fixed. The same discipline applies to inference: what is learned is
narrow. A sign's $(r,\chi)$ for a given question --- how strongly, and
in which direction, a particular bird's flight bears on ``storm'' ---
is an $E_L$ entry, learnable from labelled cases. The seated pins
$\mathrm{storm}$, $\mathrm{calm}$ are likewise $E_L$ entries, and may
be learned rather than declared. Neither requires learning a
representation in the sense of an embedding; both are learning one
small map, per sign, per question. \emph{This is the paper's one
concrete claim about the learning loop}: it is cheaper than
representation learning because it is narrow, not because it is
avoided.

\paragraph{Building a reference atom.}
\label{par:reference_atom}
A pin need not be declared; it can be built from witnessed cases the
same way live evidence is read. For a class $c$ (one storm, one
drought), let $G_c=\{y_1,\ldots,y_m\}$ be the Reading-A composite of
each witnessed episode, seated by $E_L$ from that episode's own bowl
and already rotated onto the question's axis $\hat n_q$
(\S\ref{subsec:reading_a}). Because every $y_i$ shares $\hat
n_q$ by construction, $\boxplus(G_c)$ is defined by the
companion paper's bowl-Mix equation ($\S2.7$) without needing its
still-open case of Mix across differing directions (its Future Work);
building a reference atom this way
never depends on a question that paper leaves unresolved. The
reference atom for $c$ is $\mathrm{ref}_c=\boxplus(G_c)$, its strength
reported by Equation~\eqref{eq:bijection} rather than the companion
paper's clip, per \S\ref{subsec:bijection_substitution}.

\paragraph{Coherence, and when to split.}
The capped strength $r(\mathrm{ref}_c)$ is not a safe coherence
signal on its own: it saturates at $\min(1,|z_{G_c}|)$, so a tightly
clustered class and a widely scattered one can both report the
maximal value once enough witnesses push $|z_{G_c}|$ past $1$, at
exactly the point where the underlying spread most needs
distinguishing (\S\ref{par:reference_atom_illustration} makes this
concrete). The uncapped magnitude alone does not fix this either,
since it says nothing about how much evidence produced it: $|z_{G_c}|
= 0.6$ means one thing from a single witness and another from nine
that largely cancelled.

The correct coherence signal normalizes by both at once --- the
\emph{mean resultant length}, the standard measure of angular
concentration for a set of phasors:
\begin{equation}
    \bar R_c \;=\; \frac{|z_{G_c}|}{\sum_{y\in G_c} r_y},
    \qquad
    z_{G_c} \;=\; \sum_{y\in G_c} r_y e^{i\chi_y},
    \label{eq:coherence}
\end{equation}
using the same uncapped $z_{G_c}$ the companion paper's bowl-Mix
equation already computes before capping, alongside one further sum,
$\sum r_y$, tracked once when $G_c$ is assembled and not otherwise
needed for reading live evidence. $\bar R_c\in[0,1]$: $1$ for
perfectly aligned witnesses, $0$ for complete cancellation or uniform
scatter, and --- unlike $r(\mathrm{ref}_c)$ --- unaffected by how many
witnesses there are or how strong each one is, so a class of nine and
a class of ninety are judged by the same scale. This is not a new
invention; $1-\bar R_c$ is the standard circular variance, and the
same quantity underlies tests (Rayleigh's, for one) of whether a set
of angles is distinguishable from uniform scatter at all --- a route
to choosing $\theta$ by a stated significance level rather than by
feel, not adopted here but available.

Given a coherence threshold $\theta\in(0,1)$ and a minimum
group size $k_{\min}$ below which a split is not attempted: if
$\bar R_c\ge\theta$, keep $\mathrm{ref}_c=\boxplus(G_c)$; otherwise, if
$|G_c|\ge k_{\min}$, partition $G_c$ into sub-groups and recurse on
each, so that a single label may resolve into several reference atoms
(a slow-onset and a flash drought, say); otherwise accept
$\mathrm{ref}_c$ as weak, or fall back to comparing incoming evidence
against $G_c$'s individual members rather than their compression. How
$G_c$ is partitioned when a split is warranted is not specified here
and is left open, as is the choice of $\theta$ itself; $\theta$ must
be bounded away from $0$, since every singleton group has $\bar
R_c=1$ trivially and a threshold with no floor never stops splitting.
Group size strictly decreases at each split, so the recursion
terminates regardless of $\theta$'s exact value.

\paragraph{This is level closure, applied to references.} A class
discovered to need splitting is a level-2 object in its own right
(companion paper, \S2.2): the split is discovered from $G_c$'s own
internal agreement, not asserted in advance, and each resulting
sub-reference is exactly as usable --- comparable, walkable, further
splittable --- as the class it came from. A domain with genuinely
few witnesses (two droughts, not nine) is not a failure of the
procedure; $k_{\min}$ says plainly that there is not yet evidence to
subdivide responsibly, and the honest output is a weaker reference or
a fall back to individual comparison, not a manufactured split.

\paragraph{Repeated splitting is a tree, and classifying against it
means descending it.} A label that splits once and then splits again
within a child produces a tree of reference atoms rooted at the
original class: $\mathrm{storm}$ at the root, perhaps $\mathrm{slow\
onset}$ and $\mathrm{flash}$ as children, each itself subject to the
same $\bar R\ge\theta$ test on its own witnesses. Comparing an
incoming atom against such a tree is naturally top-down: match against
the root's siblings first (storm, drought, rain, flood), then, within
whichever is nearest, descend to its children if it has any, and
repeat. This is the shape of a classical decision tree, but not its
mechanism: a decision tree tests a raw coordinate against a learned
threshold at each node; a node here compares a full state against a
reference atom by Equation~\eqref{eq:comparison_metric}, using
operators and a metric the algebra already supplies rather than
introducing threshold tests as a new primitive. The tree is discovered
by $\bar R$, level by level, the same way \S\ref{par:reference_atom}
discovers it, not designed in advance; the comparison rule and the
null of \S\ref{subsec:comparison_rule} and \S\ref{subsec:the_null}
apply at every node unchanged.

\paragraph{Illustration.}
\label{par:reference_atom_illustration}
Nine witnessed episodes at $r=0.6$ each, at
$\chi=8^\circ,9^\circ,9^\circ,10^\circ,10^\circ,10^\circ,11^\circ,
11^\circ,12^\circ$, give $\sum r_y=5.4$ and
$\mathrm{ref}=(1.00,\hat n_q,10.0^\circ)$: capped $r'=1.00$ and
$\bar R=0.9998$ agree --- coherent, confidently kept. The same nine
witnesses spread instead across
$\chi=-40^\circ,-20^\circ,-10^\circ,0^\circ,10^\circ,20^\circ,
30^\circ,45^\circ,60^\circ$ give the identical capped $r'=1.00$, since
$|z_{G_c}|$ still exceeds the cap --- but $\bar R=0.8689$, correctly
showing this class as less coherent than the first even though the
capped reading cannot tell them apart. Five witnesses at $\chi=0^\circ$
and four at $\chi=180^\circ$ --- two real sub-types wearing one
label --- give capped $r'=0.60$ and $\bar R=0.11$: both signal a
problem, but $\bar R$ states it as a small fraction of the evidence
actually agreeing, not merely a moderate reading. A perfectly even
split, five and five, gives $\bar R=0.00$, the same cancellation that
gives \S\ref{sec:worked_example} its abstain, now flagging the class
itself rather than a single reading of it.

\subsection{A walk-based classifier}
\label{subsec:walk_classifier}

Everything above classifies by comparing a state against a
\emph{compressed} reference built via $\boxplus$. The same question
--- what does new evidence resemble --- has a second mechanism,
using Reading B instead of Reading A, that has been available since
\S\ref{subsec:reading_b} and simply never assembled for this purpose:
walk new evidence to a landing, $x=\mathrm{land}_\pi(q)$, and compare
that landing against a stored reference by
Equation~\eqref{eq:comparison_metric} --- exactly what the companion
paper's storm section already gestures at (``an endpoint may fall
near the storm pin or the calm pin''), now made precise by the metric
\S\ref{subsec:comparison_rule} supplies rather than left as an
undefined ``near.'' Nothing here is new machinery; it reuses Reading
B and the comparison rule exactly as they already stand.

\paragraph{Why this classifier does not compress.}
What it cannot do, for now, is combine several witnessed landings
into one reference the way \S\ref{par:reference_atom} combines
several Reading-A composites into $\mathrm{ref}_c$. Every $y_i$
feeding a Reading-A reference atom shares $\hat n_q$ by construction
--- Reading A's own rotation step guarantees it
(\S\ref{subsec:reading_a}). A $T$-walk's landing carries no such
guarantee: nine witnessed episodes, walked independently through
their own sequences of actors, generally land nine different places.
Compressing those landings via $\boxplus$ would need exactly the
companion paper's still-open case of Mix across differing directions
(its Future Work), or a combination rule for rotations this algebra
does not currently define --- averaging rotations is a real, known
operation (a Karcher mean on $\mathrm{SO}(3)$, say) but importing one
is a decision for later, not assumed here.

\paragraph{The fallback is the whole mechanism, not a special case.}
A walk-based classifier therefore keeps every witnessed landing
individually and compares a new landing against each --- the same
fallback \S\ref{par:reference_atom} already names for a class too
small to split (``fall back to comparing incoming evidence against
$G_c$'s individual members''); here it is the only available
mechanism rather than a fallback from a failed compression. A single
declared pin, as in the storm example, is simply the case of one
witness kept rather than several.

\paragraph{Illustration.} Seat the question at the north pole,
$\hat n_q=(0,0,1)$. Two actors walk it: $r_{w_1}=\tfrac13$
($\alpha=60^\circ$) about the $x$-axis, then $r_{w_2}=\tfrac13$
($\alpha=60^\circ$) about the $y$-axis, landing at
$\hat n=(0.433,-0.866,0.25)$, i.e.\ $(\theta,\phi)=(75.5^\circ,
-63.4^\circ)$. Two declared pins: $\mathrm{dry}$ at
$(\theta,\phi)=(80^\circ,300^\circ)$ and $\mathrm{rain}$ at
$(\theta,\phi)=(30^\circ,90^\circ)$. The geodesic distance
(Equation~\eqref{eq:comparison_metric}, $w_r=w_\chi=0$) from the
landing to $\mathrm{dry}$ is $5.6^\circ$; to $\mathrm{rain}$,
$102.5^\circ$. The verdict is $\mathrm{dry}$, by a wide and
unambiguous margin, using nothing beyond a walk already defined and a
metric already stated.

\emph{Blocked for later, not solved here:} a genuine compressed
reference path, standing in for several witnessed walks the way
$\mathrm{ref}_c$ stands in for several Reading-A composites, needs
either the cross-direction Mix question resolved or an explicit
rotation-averaging operator adopted and justified on its own terms.
Both are left for iteration.

\section{Worked Example: The Diviner's Shells}
\label{sec:worked_example}

We work Reading A in full, since the companion paper never carries a
phasor reading past two states, and its property demonstration
(non-associativity, on an anonymous triple) is not a reading of a
bowl that means something.

Seat the question \emph{will it rain} at $\hat n_q$. Each shell, once
cast, is read by a diviner (or, later, by a learned $E_L$) as a state
$(r,\chi)$ at $\hat n_q$: $\chi=0$ favors rain, $\chi=\pi$ favors no
rain, and intermediate $\chi$ is a hedge. Five casts, each shell
carrying $r=0.4$ unless noted:

\begin{center}
\begin{tabular}{lcc}
\hline
Cast & reading & interpretation \\
\hline
4 shells, all favor rain      & $r=1.00$ at $\chi=0^\circ$   & confident rain \\
3 favor, 1 against            & $r=0.80$ at $\chi=0^\circ$   & rain, less confident \\
2 favor, 2 against            & $r=0.00$                     & \textbf{abstain} \\
4 favor, 2 hedge at $90^\circ$ & $r=1.00$ at $\chi\approx 26.6^\circ$ & rain, hedged direction \\
12 shells ($r=0.15$ each), all weakly favor    & $r=1.00$ at $\chi=0^\circ$   & confident rain \\
\hline
\end{tabular}
\end{center}

The third row is the case worth dwelling on. Two shells favor rain and
two oppose it, at equal strength; the phasors cancel exactly, and the
verdict is black --- not a coin flip, not a default to the majority
class, but the algebra's own statement that this evidence does not
decide the question. A system built on a forced classifier (softmax
over two classes, say) has no way to express this outcome; it must
output a confident-looking number even when the evidence is silent.
Reading A's null is not an afterthought fitted onto the algebra; it is
Axiom 2, doing here exactly what it does for a Mix reading of colour.

The fifth row shows the saturation the companion paper's \S2.7
discusses: twelve weak, agreeing signs pin at $r=1$ exactly as four
strong ones do, and the difference between ``moderately confident''
and ``overwhelmingly confident'' is lost. This is the place where
\S\ref{subsec:selection_rule}'s selection rule matters most: whether
those twelve readings are twelve units at one level, or repeated
samples of one unit that should have been reduced (companion paper,
Axiom 4) before ever reaching the bowl, is a question this paper does
not resolve and flags as the clearest candidate for the next.

\S\ref{subsec:bijection_substitution}'s substitution softens this
specific symptom without resolving the question just raised.
Calibrated to this domain's evidence with $K=0.2$, the fourth row
reads $r'=0.899$ and the fifth $r'=0.900$ --- genuinely
distinguishable, where the clip read both as $1.00$ --- and even a
hundred thousand such weak witnesses would read $r'=0.99999$, not
exactly $1$. What is now settled is that the two \emph{can} be told
apart; what remains open is whether twelve weak signs are honestly
twelve co-present witnesses of one thing, or repeated samples of a
single witness that should have been reduced before reaching the bowl
at all. The substitution reports strength faithfully either way and
does not decide which reading was owed.

\emph{Reading B on the same shells} is not worked here. Doing so
honestly requires an order for the casts, which a physical handful of
shells thrown at once does not supply --- exactly the case the swap
test (\S\ref{subsec:swap_test}) assigns to Reading A. A worked Reading
B example belongs to evidence that genuinely arrives in sequence, and
the companion paper's storm already provides one.

\section{Level Closure in Practice: Correcting a Path}
\label{sec:level_closure_practice}

A $T$-path records the sequence of actors applied to it, faithfully,
by definition (companion paper, Axiom 6). It has no means of knowing
whether that sequence was the right one --- level closure
(\S\ref{sec:inherited}) says precisely this: a lower unit carries no
representation of what it constitutes, and neither the atom nor the
path can see the higher-level context that would tell it it has erred.

A concrete case: a sequence of observed signs is walked in the order
observed, but a later sign (or a completed higher-level unit, such as
the rest of a sentence or the rest of a song) makes that order appear
wrong --- \emph{twinkle, little, twinkle, star} rather than
\emph{twinkle, twinkle, little, star}. Correcting this is a level-2
operation: it requires comparing the recorded path against context the
path itself does not have access to, and revising the walk
probabilistically. This paper's position is that the path must
therefore be \emph{legible} to the level above --- readable as an
ordered record, not only walkable --- so that a corrector at the next
level can inspect and revise it. Nothing in the companion paper's
axioms prevents this; nothing in them provides it either. We flag it
as the clearest point of contact between this paper's verdicts and a
learning loop operating above them, and leave the corrector itself
unspecified.

\section{Prediction}
\label{sec:prediction}

Everything above compares an atom to a reference built from atoms of
the \emph{same} concurrent identity: a $\mathrm{storm}$ reference is
built from patterns that were themselves storms, and a match says what
the evidence currently resembles. That is classification. Prediction
--- what is likely to happen, not what this already is --- needs a
different labelling convention, not a different mechanism.

\subsection{Immediate forecast: relabelling by outcome}
\label{subsec:immediate_forecast}

Fix a lead time $\Delta$. A \emph{precursor witness} for outcome $c$ is
a Reading-A composite $y_i$ of the evidence bowl observed at some time
$t$, labelled not by what it is but by what followed it: $c$ occurred
at $t+\Delta$. Gathering such witnesses gives a precursor group
$G_c^{\Delta}$, and everything in \S\ref{par:reference_atom} applies
unchanged: $\mathrm{ref}_c^{\Delta}=\boxplus(G_c^{\Delta})$, checked for
coherence by $\bar R_c$ (Equation~\eqref{eq:coherence}) and split if
$c$ conceals more than one precursor shape, exactly as
\S\ref{par:reference_atom_illustration} splits $\mathrm{drought}$.
Forecasting is then comparison against $\{\mathrm{ref}_c^{\Delta}\}$ by
the ordinary rule of Equation~\eqref{eq:comparison_metric}: the
nearest precursor reference, within threshold, says what is likely at
$t+\Delta$. The null of \S\ref{subsec:the_null} applies exactly as
before --- no close precursor match is an honest \emph{no forecast},
not a forced one, which matters most exactly where a wrong call is
costly.

Nothing here is a new equation. The only change from
\S\ref{subsec:learned_pins} is which witnesses are gathered under a
label: concurrent identity for classification, a fixed time-shifted
outcome for this. Two things mark it as a first, crude pass rather
than a finished one, named so they are not mistaken for settled
choices: $\Delta$ is fixed and single-valued, so the forecast has one
horizon and no notion of a nearer or more distant event both being
plausible; and no forecast carries a horizon-dependent confidence ---
a match found for $\Delta=10$ minutes and one for $\Delta=10$ days are
read with the same $\bar R_c$ scale, though the latter should
plausibly be trusted less for the same coherence value. Both are left
for iteration, not resolved here.

\subsection{A cheap decoder: transition counting over actors}
\label{subsec:markov_decoder}

Everything above answers ``what does this resemble.'' A different
question --- ``what comes next in an unfolding walk'' --- has a
mechanism simpler than either technique in
\S\ref{subsec:four_cells}, and it is worth stating precisely because
it needs none of the machinery either one does: no metric, no
$\boxplus$, no coherence check. It is a first-order Markov chain over
the discrete vocabulary $E_L$ already provides --- the same technique,
under the same name, used since the earliest work on statistical
language modelling.

Given a set of witnessed trajectories, each already an ordered
sequence of actors $(w_1,\ldots,w_n)$ (companion paper's
pipeline equation, \S2.2), count how often each actor is followed
by each other across every witnessed path:
\begin{equation}
    \widehat{P}(w_{k+1}=w' \mid w_k=w)
    \;=\;
    \frac{\mathrm{count}(w\to w' \text{ across all witnessed paths})}
         {\sum_{w''}\mathrm{count}(w\to w'' \text{ across all witnessed paths})}.
    \label{eq:markov}
\end{equation}
Decoding is a lookup: given the actor just visited, predict the $w'$
maximizing $\widehat P$, or sample from it. Nothing here touches
$T_w(x)$; the mechanism operates entirely on \emph{which} actor
follows \emph{which}, never on the rotation either one performs. It
is closer to bookkeeping on the tape a walk already produces than to
a use of the walk's own geometry.

\paragraph{Illustration.} Four witnessed paths ---
$(\text{twinkle},\text{twinkle},\text{little},\text{star})$ twice,
$(\text{twinkle},\text{little},\text{star})$ once, and
$(\text{twinkle},\text{twinkle},\text{little},\text{bat})$ once --- give
\[
\widehat P(\cdot\mid\text{twinkle}) =
\Bigl\{\text{twinkle}:\tfrac{3}{7},\ \text{little}:\tfrac{4}{7}\Bigr\},
\qquad
\widehat P(\cdot\mid\text{little}) =
\Bigl\{\text{star}:\tfrac{3}{4},\ \text{bat}:\tfrac{1}{4}\Bigr\}.
\]
A single witnessed path alone would have given an uninformative
$\tfrac{1}{1}$ split; the distinction above only exists because
several paths were witnessed, not one.

\paragraph{What this mechanism does not use, stated plainly.} The
atom carries $r$, $\hat n$, $\chi$; this decoder reads none of them,
only $\tau$-gated discrete identity. For this specific mechanism, a
bare token would do exactly as well --- the geometry is present but
idle. It is named here rather than left implicit, since claiming
otherwise for this particular technique would overstate what it
uses.

\paragraph{Cost, stated honestly.} It needs several witnessed
repetitions before its output means anything, and it only ever looks
back a short, fixed window (one prior actor here; an $n$-th order
chain extends the window at the cost of needing correspondingly more
witnessed data to fill it) --- it has no notion of the deeper
structure a full prefix match, \S\ref{subsec:path_predictors} below,
would catch. Between the two, this is the cheaper, unblocked
mechanism for the same job; the exemplar-based one is the deeper,
still-open alternative.

\subsection{Open: path predictors}
\label{subsec:path_predictors}

The immediate forecast reads only the endpoint of what happened after
a precursor; a deeper forecast reads the whole unfolding. Store, per
witnessed episode, the entire trajectory that followed --- not one
composite but the ordered sequence of atoms, timestamped as the
companion paper's step-metadata generalization allows (\S3.3) --- and
match an
\emph{observed prefix} of a trajectory now against the opening steps of
stored trajectories, reading the continuation of the closest match as
the prediction. This is not a new idea invented for this paper: it is
the analog method already used in operational forecasting, arrived at
here from the algebra's own machinery rather than borrowed from it.

It is filed open, not attempted, because it needs two things this
paper has not built. First, storage of complete historical
trajectories rather than single compressed reference atoms, a real
cost the immediate forecast does not pay. Second, and more
substantial: Equation~\eqref{eq:comparison_metric} compares one state
to one state, and prefix matching compares a \emph{sequence} of states
to a sequence, which needs a distance over sequences --- some
alignment or step-wise accumulation of the single-state metric,
robust to two episodes unfolding at different paces --- that has not
been defined. This section marks where that work belongs rather than
attempting it; the immediate forecast of
\S\ref{subsec:immediate_forecast} does not depend on it and should not
wait for it.

\subsection{Four cells, two techniques}
\label{subsec:four_cells}

Comparison (\S\ref{subsec:comparison_rule}) is the one operation
behind everything in this section and the last --- everything, that
is, except \S\ref{subsec:markov_decoder}'s transition count, which
answers a different question (what comes next, not what this
resembles) without comparison at all. Two established
techniques are built from comparison, depending on what a produced
state is checked against: \emph{nearest-centroid classification},
comparing against a single reference compressed from many witnesses
via $\boxplus$ (\S\ref{par:reference_atom}), and
\emph{nearest-neighbor classification}, comparing against witnesses
kept individually (\S\ref{subsec:walk_classifier}). Neither is
specific to this paper --- both are standard, decades-old techniques
in pattern recognition, arrived at here from the algebra's own
constraints rather than imported by name.

A second, independent axis is what the label means: identity, giving
a verdict about now, or a time-shifted outcome, giving a forecast
about later (\S\ref{subsec:immediate_forecast}). The two axes cross
to give four cells, not a new technique for each:

\begin{center}
\begin{tabular}{lll}
\hline
 & now (identity label) & later (outcome label) \\
\hline
nearest-centroid &
  \S\ref{par:reference_atom} &
  \S\ref{subsec:immediate_forecast} \\
nearest-neighbor &
  \S\ref{subsec:walk_classifier} &
  \S\ref{subsec:path_predictors} \\
\hline
\end{tabular}
\end{center}

The bottom-right cell already has its own established name for the
same reason the top row does: nearest-neighbor comparison, run
against outcome labels, over whole sequences rather than single
states, is the analog method \S\ref{subsec:path_predictors} cites.

The split between rows is not stylistic; it is forced by the same
fault line as $\boxplus$ and $T$. Compression into a centroid is
available wherever $\boxplus$'s composites are guaranteed to share a
direction; keeping individual neighbors is what remains wherever that
guarantee fails --- a $T$-landing (\S\ref{subsec:walk_classifier}),
or, in the one open cell, a whole sequence with no defined distance
yet.

Three of the four cells are specified and checked on worked numbers;
none should be read as more finished than that. Nearest-centroid
classification (\S\ref{par:reference_atom}) is complete as a
mechanism, with three choices left named and open within it: the
coherence threshold $\theta$, the comparison weights
$w_r,w_{\hat n},w_\chi$, and how a class is partitioned once a split
is warranted. Nearest-neighbor classification
(\S\ref{subsec:walk_classifier}) is complete for what it does ---
walking new evidence and comparing the landing against kept
witnesses --- with compression into a single centroid blocked
separately, not left unfinished. The immediate forecast
(\S\ref{subsec:immediate_forecast}) is explicitly a first pass: one
fixed lead time $\Delta$, no horizon-dependent confidence, both named
for iteration rather than resolved. The fourth cell,
\S\ref{subsec:path_predictors}, is open outright, waiting on the
sequence-level distance neither this table nor that section
supplies. None of the three built cells has been evaluated against
real data; each is verified on the toy numbers of its own
illustration, per the framework-versus-benchmarked choice
\S\ref{sec:scope} leaves open.

\section{Assembly: A Second Structural Verb (Status Report)}
\label{sec:assembly_status}

\textbf{Read this paragraph before anything else in this section.}
Everything below was worked out and numerically verified in
conversation, not through the same drafting and review process as
every section preceding it. None of it has been held to the citation
discipline, the cross-checking, or the settling time the rest of this
document received before being written down. It is recorded here
specifically so its state is legible --- what is closed, what is
open, and what is only a named gap with no attempt yet --- not
because any of it is ready to stand as a finished part of Verdicts.
Nothing in this section should be cited or relied upon as if it
carried the same weight as \S\ref{subsec:comparison_rule} through
\S\ref{subsec:markov_decoder}.

\subsection{What is settled: the two-body case}
\label{subsec:assembly_settled}
A second technique, distinct from comparison, was explored: two
atoms combining not by fusion ($\boxplus$) and not by bare
co-presence ($\oplus$), but by a relation added while both remain
intact.
\begin{itemize}
\item \textbf{Signature.} $A\ \mathrm{Assembly}\ B =
      (A,\,B,\,\kappa_{AB},\,U_{AB})$ --- both atoms preserved, a
      contact point and an interaction energy added.
\item \textbf{Energy, closed form.}
      $U_{AB} = -\tfrac{r_A r_B}{2}(3+\hat n_A\cdot\hat n_B)$, a
      single dot product, unconditionally defined at every angle
      including exactly $180^\circ$. Verified against an
      axis-based derivation to floating-point precision
      ($10^{-15}$) over 20{,}000 random pairs.
\item \textbf{Contact point.} $\kappa_{AB}=\hat n_A+\hat n_B$, the
      raw, unnormalized sum --- always defined, its own magnitude a
      free confidence signal.
\item \textbf{The antipodal case.} Resolved, not patched: the
      energy formula above never needed an axis at all, so the
      degeneracy that motivated days of exploration (a bisector
      dividing by zero) turned out to be an artifact of an
      unnecessary normalization step, not a property of the
      underlying relationship.
\item \textbf{Repulsion is impossible for an isolated pair.} Proven,
      not assumed: checked against 100{,}000 random pairs of varying
      orientation and strength; the least attractive result found
      was still negative.
\item \textbf{Group energy.} $\sum_{\text{pairs}} U_{ij}$, the
      standard physical form for total multi-body interaction
      energy --- computable today for any bowl.
\item \textbf{Group reconciliation, closed.} The original worry ---
      an atom cannot simultaneously occupy several private,
      independently-optimal axes with different neighbors ---
      dissolves once real position exists. With position, no axis
      is being optimized per pair at all; each pairwise force is
      simply read off real geometry, and total force is their
      ordinary vector sum, $F_{\text{total}}=\sum_i F_i$. Verified
      directly on a three-atom case: no ambiguity, no conflict,
      nothing left to reconcile beyond arithmetic already available.
\end{itemize}

\subsection{Position and distance}
\label{subsec:assembly_position}
Each atom may carry an ordinary position vector $\vec P_A$, measured
from one shared origin --- deliberately external to the bounded
$r\le 1$ stage, the same way time already sits outside the atom on
the path rather than inside it.
\begin{itemize}
\item \textbf{Distance.} $d=\lvert \vec P_B-\vec P_A\rvert$, plain
      Euclidean distance.
\item \textbf{The real axis supersedes the derived one.} Once real
      positions exist, the separation axis is
      $(\vec P_B-\vec P_A)/d$, computed from position, not derived
      from orientation. Checked directly: two atoms left exactly
      antipodal in orientation, given distinct positions, produce a
      perfectly well-defined axis and energy --- the antipodal
      problem was never triggered in the first place once position
      exists. The earlier orientation-derived axis
      (\S\ref{subsec:assembly_settled}) remains useful only as a
      fallback when position is unknown.
\item \textbf{Falloff.} The standard $1/d^3$ scaling was used
      throughout; not derived here, borrowed directly from the
      physical dipole formula this whole construction has tracked.
\item \textbf{Dynamics, closed.} Motion follows Axiom~V4: gradient
      descent, no volume term, no momentum. Verified stable at
      $N=2$, $3$, and $8$, and confirmed as the deliberate, final
      rule rather than one candidate among several --- volume-
      weighted motion was tested as the alternative and shown to
      destabilize precisely where volume ratios grow extreme, which
      settled the choice on stability grounds rather than
      convenience.
\item \textbf{Symmetric or actor/patient, closed at criterion,
      complete by design.} Polarity (\S\ref{subsec:assembly_features})
      supplies the criterion --- equal strength suggests a mutual
      bond, unequal strength suggests an actor and a patient,
      mirroring $T$'s own asymmetry. No further behavioral mechanism
      is missing: it was tested directly and found to require
      exactly the volume-weighted motion Axiom~V4 excludes on
      stability grounds. The honest, final scope of this question is
      a label describing a bond, not a difference in how it moves.
\item \textbf{Identical-atom multiset, closed.} $A$ assembling with
      an exact copy of itself needed a count, not a set --- Axiom~V5
      supplies it: $\{A,A\}$ collapses under ordinary set equality,
      $m(A)=2$ does not. Two genuinely separate witnesses sharing an
      identical tuple are now distinguished by multiplicity rather
      than lost to it, without any change to Axiom 4 or $\oplus$
      itself.
\end{itemize}

\subsection{What remains open}
\label{subsec:assembly_open}
\begin{itemize}
\item \textbf{Mass segregation (the astrophysical phenomenon),
      superseded as a prediction here.} Named early on as an
      expected consequence while this document's own dynamics was
      still built around a mass-like quantity: stronger members
      sinking toward a group's center, weaker members toward the
      periphery, as in real globular-cluster relaxation. Once
      Axiom~V4 settled motion as excluding $\mathcal{V}_{\text
      {atom}}$ entirely, this prediction's own mechanism no longer
      applies here --- confirmed directly, equal force with nothing
      dividing it gives identical movement regardless of strength or
      volume. The astrophysical phenomenon keeps its real name;
      nothing in this framework currently reproduces it, since the
      quantity that would need to drive it was deliberately kept out
      of motion. Recorded as a real idea that depended on a choice
      later reversed, not as something still to be verified.
\item \textbf{Multi-vector generalization, largely closed.} Every
      atom keeps its $N$ informational vectors, each capped at
      $r\le 1$ as usual and counted toward atomic number; alongside
      them, in the same closure unit, sits one derived vector ---
      $\sum$ of the informational vectors, uncapped, excluded from
      the count, needing no $\tau$-justified admissibility of its
      own, the same status as $\bar R_c$ or $z_{G_c}$. Coarse-grained
      assembly (the common case) runs on the derived vector, so
      $U_{AB}$ and $\kappa_{AB}$ never actually needed to generalize
      at all. What remains optional: fine-grained, simultaneous
      multi-partner bonding via the informational vectors directly,
      with two sub-mechanisms identified --- an unconstrained double
      sum over every vector pair, or an exclusive matching (the
      assignment problem, solved via the Hungarian algorithm).
      Checked precisely where these agree: sorting each atom's
      vectors by strength and pairing by rank matches the true
      optimal matching whenever alignment quality is roughly uniform
      across candidates (a consequence of the rearrangement
      inequality) but can be strictly worse than optimal when
      strength and alignment pull in different directions for
      different pairs --- verified with a constructed adversarial
      case ($-1.17$ for rank order versus $-1.68$ true optimum). The
      honestly correct general statement: the lowest-energy pairing
      is the stablest, and other pairings are unstable and will
      dynamically move toward it; strongest-with-strongest is that
      stable destination only in the uniform-alignment special case,
      not in general.
\end{itemize}


\subsection{Derived features}
\label{subsec:assembly_features}
Several quantities were derived from the existing atom, aggregate,
or pair, following the discipline already established for $\bar
R_c$ (Equation~\eqref{eq:coherence}): computed, not assigned, and
checked against what gap, if any, they close.

\begin{center}
\begin{tabular}{lp{6.5cm}l}
\hline
Feature & What it is & Status \\
\hline
Atomic number &
  Vector count per atom, proposed as bonding capacity &
  Speculative; not adopted \\
Density &
  Strength divided by vector count &
  Speculative; depends on atomic number \\
Energy ($\sum r_i$) &
  Already established
  (\S\ref{par:reference_atom}, \S\ref{subsec:bijection_substitution}) &
  Settled, predates this section \\
$z_{G_c}$, the vector sum &
  Raw, uncapped $\sum r_y e^{i\chi_y}$; direction plus absolute
  magnitude, neither carried by $\sum r_i$ or $\bar R_c$ alone &
  Verified; closes two real gaps (direction, absolute weight) \\
Polarity ($\lvert r_A-r_B\rvert$) &
  Whether a pairwise bond is balanced or one-sided &
  Verified concept; supplies the criterion for the
  symmetric/actor-patient question above \\
Heterogeneity &
  Variance of strengths within a group, distinct from $\bar R_c$'s
  directional coherence &
  Named, not built at the bowl level; the atom-level counterpart is
  now $\sigma_{\text{atom}}$ below \\
Volume concentration ($\sigma_{\text{atom}}$) &
  How evenly an atom's volume is spread across its own
  informational vectors (Equation~\eqref{eq:volume_concentration});
  $0$ even, $1$ fully concentrated, independent of $N$ &
  Verified; adopted in Axiom~V1 \\
Information energy ($\mathcal{E}_{\text{atom}}$) &
  Combined squared magnitude of an atom's own informational vectors
  as if mutually orthogonal (Equation~\eqref{eq:information_energy});
  defined directly, not as strength squared --- the same relation
  variance has to standard deviation. Distinct from bowl energy,
  which sums across witnesses rather than within one atom &
  Verified; adopted in Axiom~V1 \\
Shell concentration ($\sigma_{\text{shell}}$) &
  Whether many vectors independently cluster near the same
  magnitude (Equation~\eqref{eq:shell_concentration}); a blind spot
  of $\sigma_{\text{atom}}$, confirmed directly --- a clumped
  population read as more even than a uniform spread until this was
  built. Needs a stated shell width $w$ &
  Verified; adopted in Axiom~V1 \\
Internal alignment ($\alpha_{\text{atom}}$) &
  Strength-weighted mean cosine similarity across an atom's own
  informational vectors (Equation~\eqref{eq:internal_alignment});
  signed, so it separates active cancellation from mere
  non-reinforcement &
  Verified; adopted in Axiom~V1 \\
True angular dispersion &
  Distinguishes genuine scatter from a tight two-cluster split,
  which $\bar R_c$ and $z_{G_c}$ both read as zero either way &
  Named, not built; confirmed as a real, surviving blind spot ---
  tested against $\alpha_{\text{atom}}$ and not closed by it either
  ($-0.43$ scattered versus $-0.50$ bimodal, too near to separate) \\
Path turning-angle variance &
  How smooth or erratic a walked sequence is &
  Named, not built \\
\hline
\end{tabular}
\end{center}

Checked cross-connections: polarity closes the criterion half of the
symmetric/actor-patient question above; $z_{G_c}$'s mechanism and
true angular dispersion are plausible, untested leads toward group
reconciliation; density has a conditional connection to group
reconciliation if read as implying physical size. Atomic number does
not close a gap --- adopting it creates one (multi-vector
generalization). Energy, heterogeneity, and path turning-angle
variance were checked and found to close none of the open assembly
items.

\section{Higher Enclosure: Axioms for Position, Motion, and
  Multi-Vector Atoms}
\label{sec:higher_enclosure}

\textbf{Status note.} The four axioms below organize
\S\ref{sec:assembly_status}'s findings into the same axiomatic form
the companion paper uses, so the accumulated work has a durable,
citable shape rather than remaining a scattered set of verified
claims. They are Verdicts-native --- numbered V1--V4, deliberately
not continuing the companion paper's own Axiom 0--5, and nothing
here is proposed for that document. Each has been checked
numerically, several extensively, within this single sustained
session; none has had the multi-session settling time every axiom in
the companion paper received before being written down. Two items
named in \S\ref{subsec:assembly_open} --- the identical-atom
multiset, and the full behavioral realization of the actor/patient
question --- remain genuinely open and are not resolved by anything
below.

\subsection{Axiom V1: Informational and derived vectors}
\label{axiom:v1}
An atom's informational vectors are $N$ ordinary units of
Axiom~1 (companion paper), each independently capped at $r\le 1$;
$N$ is the atom's info number. Alongside them, in the same closure
unit, at most one \emph{derived} vector may be present:
\[
V_{\text{derived}} \;=\; \sum_{i=1}^{N} v_i,
\]
the raw, unnormalized sum of the informational vectors, computed
after and from them, never including itself. $V_{\text{derived}}$ is
uncapped and is excluded from $N$: it requires no
$\tau$-justified admissibility of its own, the same status already
held by $\bar R_c$ and $z_{G_c}$ (companion paper,
\S\ref{par:reference_atom}, \S\ref{subsec:four_cells}). Its
magnitude and direction are read as an effective $(r,\hat n)$ pair
--- $r=\lvert V_{\text{derived}}\rvert$, $\hat n =
V_{\text{derived}}/r$ --- standing in for the atom wherever a single
vector is required. For $N=1$, $V_{\text{derived}}$ coincides with
the atom's own vector and nothing changes.

Several scalar metrics follow directly from the same $N$
informational vectors, each answering a different question and none
used by any other axiom below unless explicitly named:
\[
S_{\text{atom}} = \sum_{i=1}^N r_i, \qquad
\mathcal{V}_{\text{atom}} = \sum_{i=1}^N r_i^3, \qquad
\gamma_{\text{atom}} = \frac{\mathcal{V}_{\text{atom}}}{N}.
\]
$S_{\text{atom}}$, atom strength, is the correct denominator for
atom-level coherence, $\lvert V_{\text{derived}}\rvert /
S_{\text{atom}}$, the structural twin of $\bar R_c$ one level down
--- verified to read $1.0$ for perfectly aligned informational
vectors and fall smoothly toward $0$ as they scatter.
$\mathcal{V}_{\text{atom}}$, atom volume, is strictly smaller than
$S_{\text{atom}}$ whenever any $r_i<1$, a forced consequence of the
$r\le 1$ cap rather than an assumption: cubing a number below $1$
always shrinks it, exactly as treating $r$ as a literal radius and
$r^3$ as the volume it encloses would predict. $\gamma_{\text{atom}}$,
mean vector volume, is the average per-vector share of that total.
None of the three is used in Axiom~V4; atom volume, in particular,
is deliberately excluded from motion (Axiom~V4) and exists here
purely as a descriptive metric of the atom. Earlier drafts of this
axiom called $\mathcal{V}_{\text{atom}}$ ``mass,'' $\gamma_{\text
{atom}}$ ``density,'' and $S_{\text{atom}}$ ``atom energy'' ---
all three retired as misleading: the first two because neither
participates in momentum, kinetic energy, or any calculation the
word ``mass'' would ordinarily license, and the third because a
genuine notion of energy exists here and is not $S_{\text{atom}}$.

A fourth quantity is not derived from strength by squaring it, but
stands on its own, the way variance is defined directly rather than
as squared standard deviation. Writing each informational vector as
$v_i=r_i\hat n_i$:
\begin{equation}
    \mathcal{E}_{\text{atom}} \;=\; \sum_{i=1}^N r_i^2.
    \label{eq:information_energy}
\end{equation}
$\mathcal{E}_{\text{atom}}$ is the combined squared magnitude the
atom's informational vectors would carry if every pair were mutually
orthogonal --- confirmed by direct construction, computing
$\mathcal{E}_{\text{atom}}$ from a set of strengths matches
$\lvert V_{\text{derived}}\rvert^2$ computed from those same
strengths placed at mutually orthogonal directions, exactly. This is
forced, not chosen: an exact identity places $\mathcal{E}_{\text
{atom}}$ inside the derived vector's own squared magnitude,
\[
\lvert V_{\text{derived}}\rvert^2 \;=\; \mathcal{E}_{\text{atom}}
\;+\; 2\sum_{i<j} r_ir_j(\hat n_i\cdot\hat n_j),
\]
verified to six decimal places. The structure matches the variance
of a sum, $\operatorname{Var}(\sum X_i) = \sum\operatorname{Var}
(X_i) + 2\sum_{i<j}\operatorname{Cov}(X_i,X_j)$, term for term:
$r_i^2$ stands where $\operatorname{Var}(X_i)$ stands, each vector's
own self-contained contribution; $r_ir_j(\hat n_i\cdot\hat n_j)$
stands where $\operatorname{Cov}(X_i,X_j)$ stands, how two vectors
relate. $\mathcal{E}_{\text{atom}}$ is the part that remains when
every cross term vanishes --- the atom's intrinsic energy, prior to
any question of how its own vectors align. This is distinct from
bowl energy ($\sum r_i$ across many separate witnesses,
\S\ref{par:reference_atom}): bowl energy is a sum across atoms,
$\mathcal{E}_{\text{atom}}$ a sum of squares within one.

A fifth metric answers a question the other four cannot. Atom
coherence above is a ratio of magnitudes and so cannot fall below
zero: it reads $0$ both for informational vectors that merely fail
to reinforce and for vectors that actively cancel. The
strength-weighted mean cosine similarity across all internal pairs
separates these:
\begin{equation}
    \alpha_{\text{atom}} \;=\;
    \frac{\sum_{i<j} r_i r_j\,(\hat n_i\cdot\hat n_j)}
         {\sum_{i<j} r_i r_j}.
    \label{eq:internal_alignment}
\end{equation}
Each $\hat n_i\cdot\hat n_j$ is already a cosine similarity, since
informational vectors are unit-directed; the denominator makes the
aggregate comparable across atoms of differing $N$ and differing
strengths, which the unnormalized sum is not --- two atoms with
near-identical internal arrangement but $N=2$ and $N=5$ give raw
sums of $+0.72$ and $+0.90$ against normalized values of $+0.995$
and $+0.998$. $\alpha_{\text{atom}}=+1$ exactly when all
informational vectors align; it is negative when they genuinely
oppose, reading $-1/3$ for two opposed pairs. Its floor is not
$-1$ in general: three or more vectors in three dimensions cannot
be mutually opposed, so achievable disagreement tightens as $N$
grows while achievable agreement does not --- the unnormalized sum
has maximum $\sum_{i<j}r_ir_j$, growing as $N^2$, against a hard
minimum of $-\tfrac{1}{2}\mathcal{E}_{\text{atom}}$, growing only as
$N$, which follows directly from $\lvert V_{\text{derived}}\rvert^2
= \mathcal{E}_{\text{atom}} + 2\sum_{i<j} r_ir_j(\hat n_i\cdot\hat
n_j) \ge 0$.
Agreement compounds; disagreement saturates.
$\alpha_{\text{atom}}$ is read alongside atom coherence, not in
place of it: the first asks whether the vectors agree, the second
how much of their strength survives into one resultant.

A sixth metric does for magnitude what $\alpha_{\text{atom}}$ does
for direction. $S_{\text{atom}}$ and $\mathcal{V}_{\text{atom}}$
both report a total and are silent on how that total is spread: two
atoms can carry near-identical strength ($2.200$ against $2.150$ in
a checked case) while one holds it evenly across its vectors and the
other concentrates most of it in two. Writing $s_i = r_i^3 /
\mathcal{V}_{\text{atom}}$ for the share of volume held by vector
$i$:
\begin{equation}
    \sigma_{\text{atom}} \;=\;
    \frac{N^2}{N-1}\cdot\operatorname{Var}(s_1,\ldots,s_N),
    \qquad N\ge 2,
    \label{eq:volume_concentration}
\end{equation}
with $\sigma_{\text{atom}}=0$ by convention at $N=1$.
$\sigma_{\text{atom}}=0$ exactly when volume is spread evenly and
$1$ when a single vector holds all of it, at every $N$ --- verified
at $N=2,3,4,6,8$ for both extremes. The normalization is what earns
this form over the simpler candidates: the largest vector's share
and the Herfindahl index were both tested and rejected, since each
reads a perfectly even atom as $1/N$ and so slides from $0.5$ to
$0.125$ as $N$ grows from $2$ to $8$, confounding evenness with
size. $\mathcal{V}_{\text{atom}}$ rather than $S_{\text{atom}}$
supplies the shares deliberately: cubing amplifies concentration, so
a single dominant vector at $r=0.95$ against three at $r=0.2$
registers a share of $0.973$, where linear strength would report the
same atom as unremarkable.

A seventh metric answers a question none of the first six can: not
how strength or volume is shared among distinct vectors, but whether
many vectors independently arrive at nearly the same magnitude.
$\sigma_{\text{atom}}$ is blind to this by construction --- four
hundred vectors at an identical $r=0.65$ each receive their own tiny
individual share, indistinguishable to
Equation~\eqref{eq:volume_concentration} from four hundred vectors
scattered across the full range; checked directly, a genuinely
clumped population registered as \emph{more} even
($\sigma_{\text{atom}}=0.00032$) than a uniform spread of the same
size ($0.00214$). Fixing this needs a shell width $w$, a named,
configurable parameter rather than a derived one, defaulting to
$w=0.1$ --- ten shells spanning $[0,1]$. Binning the $N$
informational vectors by $\lfloor r_i/w\rfloor$ into $K$ occupied
shells with populations $n_1,\ldots,n_K$ and $p_k = n_k/N$:
\begin{equation}
    \sigma_{\text{shell}}(w) \;=\;
    \begin{cases}
        1, & K=1 \\[4pt]
        \dfrac{K^2}{K-1}\cdot\operatorname{Var}(p_1,\ldots,p_K), & K\ge 2.
    \end{cases}
    \label{eq:shell_concentration}
\end{equation}
The $K=1$ case is stated separately because the general form has a
removable singularity there: every vector sharing one shell is
maximal concentration, not the zero the formula would otherwise
return. $\sigma_{\text{shell}}=0$ when occupied shells are evenly
populated, at every $K$ tested; $\sigma_{\text{shell}}=1$ when every
vector shares one shell. Checked at $w=0.1$ and $w=0.05$: well
separated shells (as in $10@0.98,\,3@0.3,\,400@0.65$, giving
$\sigma_{\text{shell}}=0.908$ at either width) are stable to the
choice, while genuinely uniform populations are not --- occupied
shell count nearly doubled between the two widths in a checked
uniform case, though the concentration score itself stayed
negligible either way. $w$ should be reported alongside
$\sigma_{\text{shell}}$ whenever it is used.

\subsection{Axiom V2: External position}
\label{axiom:v2}
An atom, or a bowl, may carry a position vector $\vec P$, measured
from one shared origin, in ordinary unbounded Euclidean space ---
deliberately external to the $r\le 1$ stage, as time already sits
external to the atom on the path rather than inside it. For two
atoms $A,B$ with positions $\vec P_A,\vec P_B$:
\[
d_{AB} = \lvert \vec P_B - \vec P_A\rvert, \qquad
\hat r_{AB} = \frac{\vec P_B-\vec P_A}{d_{AB}}.
\]
Where position exists, $\hat r_{AB}$ supersedes any
orientation-derived separation axis; the latter remains a fallback
only where position is unknown.

\subsection{Axiom V3: Assembly}
\label{axiom:v3}
Two atoms combine by a relation added while both remain intact ---
distinct from $\oplus$ (bare co-presence, no relation) and $\boxplus$
(fusion, identities lost):
\[
A\ \mathrm{Assembly}\ B \;=\; (A,\,B,\,\kappa_{AB},\,U_{AB}),
\]
\[
U_{AB} \;=\; -\frac{r_A r_B}{2}\cdot
\frac{3+\hat n_A\cdot\hat n_B}{(d_{AB}^2+\varepsilon_{AB}^2)^{3/2}},
\qquad
\kappa_{AB} = \hat n_A+\hat n_B,
\qquad
\varepsilon_{AB} = r_A+r_B.
\]
$U_{AB}$ is unconditionally defined at every angle, including exactly
antipodal orientation, and at every separation, including $d_{AB}=0$.
For an isolated pair, $U_{AB}<0$ always: repulsion does not occur
between two atoms considered alone (verified over 100{,}000 random
pairs). For a group, total energy is $\sum_{\text{pairs}}U_{ij}$;
reconciling several simultaneous relationships on one atom requires
no separate mechanism once Axiom~V2 holds, since each pairwise force
is read from real position rather than independently optimized.
$r_A,r_B$ in the above are each atom's $(r,\hat n)$ under
Axiom~V1 --- an atom's own vector if $N=1$, its derived vector
otherwise.

\subsection{Axiom V4: Motion without atom volume}
\label{axiom:v4}
Position changes by moving directly along the energy gradient. No
term for $\mathcal{V}_{\text{atom}}$ appears, by deliberate design
rather than by an unresolved choice: atom volume (Axiom~V1) is a
real, useful metric of an atom's own structure, but makes no sense
as an input to motion, and is excluded from it.
\[
\vec P_i \;\leftarrow\; \vec P_i + dt\cdot F_i, \qquad
F_i = -\nabla_{\vec P_i}\sum_{j\ne i} U_{ij}.
\]
No inertia means no overshoot: an isolated pair closes distance
monotonically and never exceeds its own starting separation
(verified). A group of mutually attracting atoms converges to a
stable configuration centered on
$\vec P_{\text{center}}=\sum_i r_i\vec P_i / \sum_i r_i$, held
stationary once reached rather than merely passed through ---
verified at $N=2$, $N=3$, and $N=8$, flat to four decimal places
over thousands of subsequent steps in each case. $\varepsilon_{AB}$
(Axiom~V3) must be sized relative to $dt$ or the same instability a
literal $1/d^3$ term would produce at close approach can recur.

A direct, checked consequence: since $F_A=-F_B$ exactly for any
pair, and nothing divides the resulting step, two atoms of any
polarity move by identical magnitude toward each other, regardless
of strength. Weighting each atom's step by its own
$\mathcal{V}_{\text{atom}}$ was tested directly as an alternative
and produces the intended differential movement, but destabilizes
into oscillation precisely where that ratio is most extreme ---
confirmed at a $9\times$ strength difference, giving a $729\times$
volume ratio and a demonstrably non-monotonic trajectory, against
the clean monotonic decline of this axiom's rule under the same
pair. Axiom~V4 is therefore stated without $\mathcal{V}_{\text
{atom}}$ on stability grounds, not merely convention; a
differential-motion mechanism is not a matter this axiom leaves
open, since any such mechanism reintroduces the instability
directly.

\subsection{Axiom V5: Multiplicity}
\label{axiom:v5}
A bowl's membership under the companion paper's Axiom 4 is binary:
an atom is present or absent, since $\oplus$ is defined as ordinary
set union on a subset of $\mathcal{S}$. Where the number of times a
value was independently witnessed must itself survive --- not
merely whether it occurred --- a bowl may additionally be read
through a multiplicity function:
\[
m:\mathcal{S}\to\mathbb{N}_0, \qquad
m(A) = \text{the number of independent witnesses producing exactly } A.
\]
An ordinary bowl is the special case $m(A)\in\{0,1\}$ for every $A$,
recovering the companion paper's Axiom 4 exactly; nothing about
$\oplus$ or Axiom 4 is altered, and $m$ is read alongside a bowl,
never in place of it. Where $m(A)\ge 2$, $A$ is counted that many
times in every aggregate drawn from the bowl --- $\sum r_i$, $\bar
R_c$ (Equation~\eqref{eq:coherence}), and any use of Assembly
(Axiom~V3) that draws on the bowl's members individually. Two atoms
sharing an identical $(r,\hat n,\chi,\tau)$ tuple, arising from
genuinely separate witnesses, are distinguished by their
multiplicity alone, never by any difference in the tuple itself:
$A$ assembling with an exact copy of itself is simply $m(A)=2$,
unremarkable once presence and count are kept as separate
questions. Treating such a pair as the set $\{A,A\}$ is not a
looser way of saying the same thing; $\{A,A\}=\{A\}$ under ordinary
set equality, and a second, genuinely distinct witness is silently
lost --- confirmed directly: an atom of strength $r=0.7$, witnessed
twice, contributes $0.7$ to $\sum r_i$ if read as a set and the
correct $1.4$ if read through $m$.

\section{Scope and Open Questions}
\label{sec:scope}

This paper is a framework, in the same sense the companion paper is:
it specifies mechanism and states honestly what it has not
demonstrated at scale. \emph{[OPEN --- decide before circulating: does
this paper stop here, as a specification plus one worked toy example,
or does it commit to an evaluation against real data (e.g.\ a labelled
forecasting or forex task) with a stated baseline? The two choices
imply different remaining sections and different claims in the
abstract; see the companion paper's own title history for how much
that choice shapes what a title may honestly claim.]}

Regardless of that choice, the following are stated open rather than
resolved by this paper:

\begin{itemize}
\item The comparison rule (\S\ref{subsec:comparison_rule}) is a
      choice, not a derivation; a different metric may be more honest
      for a given domain.
\item The selection rule (\S\ref{subsec:selection_rule}), including
      how the grouping window of companion-paper \S2.2 is set, is the
      largest source of real difficulty and is entirely unaddressed
      here.
\item The saturation--under--aggregation problem
      (\S\ref{sec:worked_example}, companion paper \S2.7 and \S2.3) is
      diagnosed, not solved: whether to route repeated samples through
      a reduction before the bowl, or revisit the companion paper's
      Axiom~1 note on an order-preserving bijective bound, is left to
      whichever setting first needs to aggregate at scale.
\item The path corrector (\S\ref{sec:level_closure_practice}) is named
      and not built.
\item Path predictors (\S\ref{subsec:path_predictors}) are named and
      not built; the immediate forecast
      (\S\ref{subsec:immediate_forecast}) does not require them.
\item The coherence threshold $\theta$ and the method for
      partitioning a class once a split is warranted
      (\S\ref{par:reference_atom}) are open choices not fixed here;
      \S\ref{subsec:four_cells} consolidates this and every other item
      above by which of the four classification-and-forecast cells it
      belongs to.
\end{itemize}

\end{document}
